\documentclass[12pt,a4wide]{report}

\usepackage{amsthm,amssymb,mathrsfs,setspace,pstcol}%amsmath, latexsym,footmisc

\usepackage{play}
\usepackage{epsfig}
%\usepackage[grey,times]{quotchap}
\usepackage[nottoc]{tocbibind}
\renewcommand{\chaptermark}[1]{\markboth{#1}{}}
\renewcommand{\sectionmark}[1]{\markright{\thesection\ #1}}
%

% ---- Packages added for this report ----
\usepackage{amsmath}
\usepackage{graphicx}
\usepackage{booktabs}
\usepackage{float}
\usepackage{url}
\usepackage{tikz}
\usetikzlibrary{arrows.meta,positioning}
\graphicspath{{figures/}}

\input xy
\xyoption{all}


\theoremstyle{plain}
\newtheorem{theorem}{Theorem}[section]
\newtheorem{lemma}[theorem]{Lemma}
\newtheorem{corollary}[theorem]{Corollary}
\newtheorem{proposition}[theorem]{Proposition}

\theoremstyle{definition}
\newtheorem{definition}[theorem]{Definition}
\newtheorem{example}[theorem]{Example}
\newtheorem{notation}[theorem]{Notation}

\theoremstyle{remark}
\newtheorem{remark}[theorem]{Remark}

\renewcommand{\baselinestretch}{1.5}
\begin{document}

% --------------- Title page -----------------------
\begin{titlepage}
\enlargethispage{2cm}

\begin{center}

\vspace*{-1cm}

\textbf{\Large MALAYALAM NAMED ENTITY RECOGNITION USING FINE-TUNED MULTILINGUAL BERT}\\[10pt]

\vspace*{1.5 cm}

                         A Project Report Submitted \\
                     in Partial Fulfilment of the Requirements  \\
                     for the Course CSE413 Speech and Natural Language Processing \\
                     of the Degree of  \\
                          \vspace{4.5mm}
                   {\Large \bf Bachelor of Technology } \\
                   in \\
                   {\large \bf Computer Science and Engineering } \\

                      \vspace{9mm}
                   {\em  by} \\ \vspace{2.5mm}
             {\large \bf Albin John} {\large \bf (Roll No. 2023BCD0005)}\\[4pt]
             {\large \bf Adnan Soni} {\large \bf (Roll No. 2023BCD0017)}\\[4pt]
             {\large \bf Aswin Paramesh S S} {\large \bf (Roll No. 2023BCD0024)}\\[.35in]

\vfill

\begin{figure}[h]
  \begin{center}
  \includegraphics[width=6cm, height=3cm]{logo2.jpg}
  \end{center}
\end{figure}
\vspace*{0.2cm}
{\em\large to }\\%[8pt]
{\bf \mbox{DEPARTMENT OF {\bf COMPUTER SCIENCE AND ENGINEERING }}} \\%[4pt]
{\bf \mbox{INDIAN INSTITUTE OF INFORMATION TECHNOLOGY}}\\%[4pt]
{\bf KOTTAYAM-686635, INDIA}\\%[8pt]
{\it October 2026}

\end{center}

\end{titlepage}

\clearpage
\pagenumbering{roman} \setcounter{page}{2}

% --------------- Abstract page -----------------------
\begin{center}
{\Large{\bf{ABSTRACT}}}
\end{center}

Named Entity Recognition (NER) identifies proper names in text and
classifies them as persons, organizations or locations. It is hard for
Malayalam, a low-resource and highly agglutinative language that has no
capitalization and fuses case markers directly onto names.

Following Bharathi Mohan \emph{et al.}~\cite{mohan2023finetuned}, we fine-tune
a multilingual BERT encoder, \texttt{ai4bharat/IndicBERTv2-MLM-only}, on the
Malayalam portion of the Naamapadam corpus~\cite{mhaske2023naamapadam}
(IOB2 tags; PER, ORG, LOC). Agglutination is handled by aligning each word's
label with its first subword. We compare two output layers: an independent
softmax classifier trained with class-weighted cross-entropy, and a
Conditional Random Field (CRF) that models tag transitions.

On the manually annotated test set (965 sentences, 1{,}505 entities), the
CRF model reaches a strict entity-level F1 of \textbf{84.25\%} and a token
accuracy of 94.08\%. This is 12.0 F1 points above the softmax head (72.23\%),
with about 80\% fewer invalid tag sequences. PER (90.6 F1) and LOC (85.5 F1)
are recognised reliably, while multi-word ORG names (66.6 F1) remain the main
source of error. The model is deployed through a Gradio web interface and a
FastAPI REST service.

\vspace{0.5cm}
\noindent\textbf{Keywords:} Named Entity Recognition, Malayalam, Multilingual
BERT, IndicBERTv2, Conditional Random Field

\clearpage

\tableofcontents

\newpage

\pagenumbering{arabic}
\setcounter{page}{1}

% =========== Main chapters ===========

\input chapter1.tex

\input chapter2.tex

\input chapter3.tex

\input chapter4.tex

\input chapter5.tex

\bibliographystyle{plain}
\bibliography{bib}

\end{document}
